\par\Needspace{4\baselineskip}
\begin{center}
\hypertarget{native-057-article039}{}\pdfbookmark[2]{39.}{article-039}39.
\end{center}
It is immediately evident that nearly half the work required by those trials could be saved by arranging the table so that $\log B$ can be read off directly with $A$ as the argument. Three operations would then remain. The first is indirect and consists in determining $A$ so that equation 1) of Article 37 is satisfied. The second, determining $E$ from $A$ and $B$, is performed directly by either the equation $E=2B(A^{\frac{1}{2}}+\tfrac{1}{15}A^{\frac{3}{2}})$ or the equation $\sin E=2B(A^{\frac{1}{2}}-\tfrac{3}{5}A^{\frac{3}{2}})$. The third is to determine $v$ from $E$ by equation VII of Article 8. We shall reduce the first operation to a rapid algorithm free of unguided trials. By entering a new quantity $C$ in our table, we shall combine the second and third operations into one. This dispenses with $E$ altogether and gives us an elegant and convenient formula for the radius vector. We shall now take these points in order.

First we shall transform equation 1) so that \textsc{Barker}'s table can be used to solve it. To this end we set $A^{\frac{1}{2}}=\operatorname{tang}\tfrac{1}{2}w\,\mathord{.}\,\BareRoot{\frac{5-5e}{1+9e}}$, which gives
\[
75\operatorname{tang}\tfrac{1}{2}w+25\operatorname{tang}\tfrac{1}{2}w^3=\frac{75kt\BareRoot{(\tfrac{1}{5}+\tfrac{9}{5}e)}}{2Bq^{\frac{3}{2}}}=\frac{\alpha t}{B},
\]
denoting the constant $\frac{75k\BareRoot{(\tfrac{1}{5}+\tfrac{9}{5}e)}}{2q^{\frac{3}{2}}}$ by $\alpha$\HistoricalNoteEight{Manuscript note:] If \textsc{Burckhardt}'s table is used [Connaissance des Temps 1818, see also vol. VI, p. 596], then $\frac{\BareRoot{(\tfrac{1}{10}+\tfrac{9}{10}e)}}{q^{\frac{3}{2}}}=\alpha$.}. Thus, if $B$ were known, $w$ could be taken immediately \SourceLeaf{058}from \textsc{Barker}'s table as the true anomaly corresponding to the mean motion $\frac{\alpha t}{B}$. From $w$ we obtain $A$ by the formula $A=\beta\operatorname{tang}\tfrac{1}{2}w^2$, denoting the constant $\frac{5-5e}{1+9e}$ by $\beta$. Although $B$ only becomes known from $A$ by means of our auxiliary table, its difference from unity is so small that $w$ and $A$ can be expected to be only slightly in error if the divisor $B$ is omitted altogether at first. We shall therefore first find rough values of $w$ and $A$ by setting $B=1$. With the approximate value of $A$ we find $B$ itself from our auxiliary table and use it to repeat the calculation more accurately. In most cases the corrected value of $A$ will give exactly the same value of $B$ as the approximate one, so a further repetition will be superfluous except where the value of $E$ is already quite considerable. It scarcely needs pointing out that if an approximate value of $B$ happens to be known from any other source at the outset (as will always be the case when one or two of several positions close to one another have already been computed), it is preferable to use that value in the first approximation. Anyone practised in calculation will thus very often have no need to repeat it even once. This very rapid approximation is possible because $B$ differs from $1$ only by a quantity of the fourth order, which is moreover multiplied by a very small numerical coefficient. It should now be clear how this advantage was secured by introducing the quantities $E-\sin E$ and $\tfrac{9}{10}E+\tfrac{1}{10}\sin E$ in place of $E$ and $\sin E$.
